%% ⑤ 국문 초록 — docx 기준
%% 제목: 20pt / 학과·전공·성명·지도교수: 14pt
%% 본문: 명조체 12pt, 줄간격 160%, 들여쓰기 10pt
%% 2페이지 이내, 끝에 주제어 5단어 내외
\clearpage
\thispagestyle{fancy}
\addcontentsline{toc}{chapter}{Abstract (Korean)}

\begin{flushright}
  {\fontsize{12}{20}\selectfont (국문초록)}
\end{flushright}


\begin{center}
  {\fontsize{20}{26}\selectfont \thesistitleko}

  {\fontsize{15}{26}\selectfont \thesistitlekosub}
\end{center}

\vspace{5mm}

\begin{flushright}
  {\fontsize{14}{18}\selectfont \gradschool\quad \department}\\
  {\fontsize{14}{18}\selectfont \major}\\
  {\fontsize{14}{18}\selectfont \authorname}\\
  {\fontsize{14}{18}\selectfont 지도교수：\advisorname}
\end{flushright}

\vspace{4mm}

%% 본문: 12pt, 160% 줄간격, 들여쓰기 10pt
\begingroup
\fontsize{12}{12}\selectfont
\setlength{\parindent}{10pt}
\setlength{\parskip}{0pt}
\spaceskip=0.6em
초록(abstract)은 학위논문 전체의 개요(overview) 또는 요약(summary)에 해당되는 내용으로 구성되며, 논문의 주요 내용(main story)과 몇 가지 세부 사항을 포함한다. 논문 전체를 읽는 독자나 그렇지 않은 독자에게 모두 명확하여야 한다.

초록에 기술되어야 하는 주요 내용은 1) 문제의 정의(problem definition), 2) 문제의 답을 도출하기 위하여 사용한 방법 또는 도구, 3) 연구를 통하여 얻은 결과 및 결론, 4) 가능하다면 연구결과에 대한 의미, 파급효과 등에 대한 토의(discussion)를 포함한다.

초록은 일반적으로 하나의 문단(paragraph)으로 작성하며, 한 문단에서 내용이 흐르게 연속성을 유지하고, 흐름을 위한 신호(signal) 단어를 사용한다. (예, We found that\ldots, We concluded that\ldots, We suggest that\ldots)
\endgroup

\vspace{4mm}

\noindent{\fontsize{12}{16}\selectfont 주제어:} {\fontsize{12}{16}\selectfont 단국대학교, 대학원, 공학계열, 학위논문, 작성양식}

\newpage
